% =============================================================================
% Lecture 8b: Model Training, Part 2: Fitting the Step
% Introduction to Machine Learning Systems (Volume I), Chapter 8
% =============================================================================
% Thesis: a training step that needs seven accelerators' worth of memory fits
% on one once each memory term meets the technique that removes it, each at a
% cost the engineer must name. Running case: GPT-2 XL at batch 32 of 1,024
% tokens on one H100, 598 GB to 311 to 109.6 to 66.8 GB.
% Hook and synthesis: AlexNet's two 3 GB GTX 580s (Model Training, Sec. 8.7.1).
%
% Build: make -C slides build-08b_training   (figures: images/make_figs.py)
% Plan:  lecture_map.yml and storyboard.md in this folder.
% =============================================================================

\documentclass[aspectratio=169, 12pt]{beamer}
\usepackage{../../assets/beamerthememlsys}
\mlsyssetup{
  volume={Volume I},
  chapter={Chapter 8},
  chaptertitle={Model Training, Part 2},
  layer=training,
  logo={../../assets/img/logo-mlsysbook.png},
  instlogo={../../assets/img/logo-harvard.png},
  dividerstack,
}
\usepackage[T1]{fontenc}
\usepackage[scaled=0.9]{helvet}
\usepackage{courier}
\renewcommand{\familydefault}{\sfdefault}
\usepackage{booktabs,amsmath}
\usetikzlibrary{arrows.meta,positioning,calc}
\graphicspath{{images/}}
\setcounter{mlsystotalsections}{5}

\title{Model Training, Part 2}
\author{Vijay Janapa Reddi}
\institute{Harvard University}
\date{}

% A large answer choice for predict and concept-question slides
\newcommand{\choice}[2]{\textbf{\textcolor{mlsysaccent}{#1}}\enspace #2}
% One score cell of the online-softmax row: x, value, fill color
\newcommand{\scorecell}[3]{%
  \node[draw=darktext, fill=#3, minimum width=1.2cm, minimum height=0.75cm,
        font=\normalsize\bfseries] at (#1,0) {#2};}



\begin{document}

% =============================================================================
% TITLE
% =============================================================================
\mlsysnote{[0.5 min] [CORE] Read the thesis aloud once. Last lecture counted the bytes of one GPT-2 XL training step and found it far too large for one accelerator. Today removes those bytes one technique at a time and names what each one costs. The notes total about 71 minutes. FLEX: notes tagged [OPTIONAL] can be cut; cutting all of them leaves a complete lecture of about 51.5 minutes.}
\mlsystitle{Model Training, Part~2}{A step that needs seven accelerators' memory fits on one once each memory term meets the technique that removes it, at a cost the engineer must name.}{cover_training_blueprint.png}

% =============================================================================
% SECTION: THE STEP
% =============================================================================
\mlsysnote{[0.5 min] [CORE] Open with a model that did not fit on one accelerator in 2012, and one that does not fit today.}
\mlsysdivider{How many accelerators does one training step need?}{Model Training, \S8.6.6 and \S8.7.1}
\section{The Step}

\begin{frame}{In 2012, AlexNet needed two GPUs because one was too small}
  \begin{columns}[T]
    \begin{column}{0.44\textwidth}
      \centering\mlsysphoto[width=\textwidth,height=4.3cm,keepaspectratio]{nvidia-gtx580-die.jpg}\\[0.1cm]
      {\small\textbf{GTX 580, 2012: 3 GB.}}\\
      {\footnotesize\textcolor{midgray}{AlexNet: 60M parameters, split across two.}}
    \end{column}
    \begin{column}{0.52\textwidth}
      \centering
      \uncover<2->{\mlsysphoto[width=\textwidth,height=4.3cm,keepaspectratio]{nvidia-h100-pcie.jpg}\\[0.1cm]
      {\small\textbf{H100, 2022: 80 GB.}}\\
      {\footnotesize\textcolor{midgray}{GPT-2 XL, batch 32: 598 GB per step in FP32.}}}
    \end{column}
  \end{columns}
  \mlsyscredit{Photos: Fritzchens Fritz (CC0); Geekerwan (CC BY 3.0); Wikimedia Commons. AlexNet: Krizhevsky et al. 2012, via Model Training, \S8.7.1. GPT-2 XL: Napkin Math 8.10.}
  \note{[2 min] [CORE] Step 1: recall Lecture 1's AlexNet, 60 million parameters split across two 3 GB GTX 580 cards because one could not hold the training run; the photo is that card's die. Step 2: an H100, with 80 GB, about 27 times more. Today's running example, GPT-2 XL (1.5 billion parameters) trained in batches of 32 sequences of 1,024 tokens, needs 598 GB per step if every number is stored in 32 bits. ENGAGE: in ten years accelerator memory grew about 27 times; did the problem go away? (Expected: no; models and batches grew faster.) WARN: the die photo shows the processor chip only; the 3 GB of memory sat in separate chips on the card.}
\end{frame}

\begin{frame}{How many H100s will this step need by the end?}
  \mlsysactivity{Predict}{2 min}
  \centering
  \safeimg[width=0.80\textwidth,height=3.9cm,keepaspectratio]{scoreboard_step1.pdf}\\[0.2cm]
  {\normalsize Same model, same batch of 32, same arithmetic per update.}\\[0.2cm]
  {\large\choice{A}{8}\qquad\choice{B}{4}\qquad\choice{C}{2}\qquad\choice{D}{1}}
  \mlsyscredit{Model Training, Napkin Math 8.10 (FP32, standard attention, batch 32 of 1,024 tokens). The H100's 80 GB is 80 GiB, or 85.9 GB in the decimal units used for tensor sizes.}
  \note{[2 min] [CORE] The scoreboard from Part 1. Gray is the 24 GB of static state (weights, gradients, optimizer state); green is 574 GB of activations, the values saved for the backward pass. The dashed line is one H100: 80 GB is 80 GiB (binary gigabytes of $2^{30}$ bytes), or 85.9 GB in the decimal units the chapter uses. 598 GB is about seven H100s of memory. ENGAGE: students vote with fingers, then write their answer down. Do not reveal it; the scoreboard returns at the end of each segment. Most students choose A or B, because the arithmetic does not change and they assume memory scales with it. That intuition is today's target. WARN: the question is about memory; one H100 still takes about 35 days for the run.}
\end{frame}

\begin{frame}{Recall: what does one training step hold?}
  \mlsysactivity{Retrieval}{1.5 min}
  \centering
  {\Large\bfseries Question of the day}\\[0.2cm]
  {\normalsize Write the four kinds of state a step holds, the bytes per parameter\\of mixed-precision Adam, and which term grows with batch size.}\\[0.4cm]
  \uncover<2->{%
    \safeimg[width=0.85\textwidth,height=2.2cm,keepaspectratio]{static_bytes.pdf}\\[0.1cm]
    {\normalsize Batch 32, mixed precision: 24 static $+$ 86 activations $+$ 201 scores $=$ \textbf{311 GB}}}
  \mlsyscredit{Model Training, \S8.3.2 and Napkin Math 8.3.}
  \note{[1.5 min] [CORE] Retrieval from Part 1. Give 60 seconds. Step 2: weights, gradients, optimizer state, and activations. Mixed-precision Adam keeps 2 bytes of 16-bit weights, 2 of 16-bit gradients, 4 of a 32-bit master copy, and 8 for Adam's two 32-bit moments, 16 bytes per parameter or 24 GB for GPT-2 XL. Only the activations grow with batch size and sequence length. At batch 32 in mixed precision they are about 287 GB, of which about 201 GB are attention scores. For a student who missed Part 1, this slide is the whole starting point.}
\end{frame}

\begin{frame}{How do we fit a step that needs seven accelerators onto one?}
  \normalsize
  By the end of today, you will be able to:\\[0.2cm]
  \begin{itemize}\setlength{\itemsep}{0.22cm}
    \item Explain what FP16 and BF16 keep and lose, and why 32 bits remain.
    \item Trace where attention's bytes go, and compute softmax one block at a time.
    \item Size a micro-batch and a checkpoint schedule for a step that does not fit.
    \item Name what each technique costs: numerics, arithmetic, or update frequency.
  \end{itemize}
  \vspace{0.15cm}
  {\small\textcolor{midgray}{Running case: GPT-2 XL, 48 layers, batch 32 $\times$ 1,024 tokens, one H100 (80 GiB = 85.9 GB).}}\\[0.2cm]
  \threadstrip{0}{1}{0}{0}{1}
  \note{[1.5 min] [CORE] The title is the driving question. Each line is one segment, precision, tiling, and trading time for memory; the fourth line is the question asked at the end of every segment. The running case stays the same all lecture, GPT-2 XL with 48 transformer layers with hidden width 1,600 and 25 attention heads, trained on 32 sequences of 1,024 tokens per step on one H100 with 85.9 GB of high-bandwidth memory (HBM). LINK: Lecture 7 showed the calls (autocast, GradScaler, and torch.utils.checkpoint); today asks when and why to use each one, and what it costs. The strip fills the machine factor and the fast-enough half, because each technique changes what the machine must hold and how fast it runs while the model stays the same.}
\end{frame}

% =============================================================================
% SECTION: PRECISION
% =============================================================================
\mlsysnote{[0.5 min] [CORE] The first technique changes the number format. Start with what a format is.}
\mlsysdivider{What does a 16-bit number keep, and what does it lose?}{Model Training, \S8.6.3}
\section{Precision}

\begin{frame}{An FP32 run leaves 14 of every 15 H100 FLOPs unused}
  \begin{columns}[c]
    \begin{column}{0.68\textwidth}
      \centering\mlsyssteps[width=\textwidth,height=4.8cm,keepaspectratio]{precision_ladder_step}{3}
    \end{column}
    \begin{column}{0.29\textwidth}
      \centering\mlsysphoto[width=\textwidth,height=2.8cm,keepaspectratio]{google-tpu-v3-board.jpg}\\[0.1cm]
      {\footnotesize\textcolor{midgray}{Google TPU v3 board: TPUs pioneered BF16 for deep learning training.}}
    \end{column}
  \end{columns}
  \mlsyscredit{Model Training, Fig.\ 10 (NVIDIA datasheets, dense rates). Photo: Zinskauf (CC BY-SA 4.0), Wikimedia Commons.}
  \note{[1.5 min] [OPTIONAL] Peak dense arithmetic rate by number format, log scale. Step 1: the V100 of 2017: 15.7 TFLOP/s in FP32 on ordinary cores, 125 in 16-bit on Tensor Cores. Step 2: four generations. From V100 to B200, FP32 grew 4.8 times while 16-bit matrix rates grew 18 times; Google's TPU (right) pioneered BF16 along the same path. Step 3: the H100 delivers 67 TFLOP/s in FP32 against 989 in FP16 or BF16, 15 times more. Staying in FP32 leaves 14 of every 15 FLOPs unused. TF32 (494.5) is NVIDIA's 19-bit format for FP32 inputs; FP8 and FP4 return in Lecture 10b. WARN: NVIDIA's H100 table lists 1,979 BF16 with sparsity; dense rates are half, 989. ENGAGE: why do hardware designers favor narrow formats? (Expected: more multipliers per silicon area and fewer bytes per operand.) FLEX: cut freely.}
\end{frame}

\begin{frame}{FP16 gives up range; BF16 gives up precision}
  \centering
  \mlsyssteps[width=0.95\textwidth,height=5.2cm,keepaspectratio]{fp_formats_step}{4}
  \mlsyscredit{Model Training, \S8.6.3 and table 9. FP32 and BF16 range ends from the IEEE 754 and bfloat16 formats.}
  \note{[2.5 min] [CORE] Each row is one format, with the bits to scale on the left the range of normal values on a log axis on the right. Step 1: FP32 has 1 sign bit, 8 exponent bits, and 23 mantissa bits. The exponent sets the range; the mantissa sets the digits kept, about seven decimal digits. Step 2: FP16 keeps 5 exponent bits and 10 mantissa bits. Its normal range shrinks to $6.1\times10^{-5}$ through 65,504, about three digits. The pale band is the subnormal range, tiny values below the normal floor, down to about $6\times10^{-8}$, kept with fewer digits; many accelerators flush them to zero. Step 3: BF16 keeps FP32's 8 exponent bits and cuts the mantissa to 7: FP32's range, two to three digits. Step 4: a gradient of $10^{-6}$ is an ordinary number in FP32 and BF16 and falls below FP16's normal floor. LINK: Lecture 5 showed gradients shrinking deep in the backward pass. ENGAGE: which format would you pick for gradients, and why? (Expected: BF16, for range.)}
\end{frame}

\begin{frame}{A small update vanishes into a 16-bit weight}
  \mlsysactivity{Predict}{2 min}
  \centering
  {\normalsize A weight of 0.5 is stored in FP16. The optimizer adds 0.0001.\\What is stored now? After 10,000 such updates?}\\[0.3cm]
  \uncover<2->{%
    \only<-2|handout:0>{\safeimg[width=0.92\textwidth,height=4.0cm,keepaspectratio]{swamp_step1.pdf}}%
    \only<3-|handout:1>{\safeimg[width=0.92\textwidth,height=4.0cm,keepaspectratio]{swamp_step2.pdf}}}
  \mlsyscredit{Derived from the IEEE 754 FP16 format; checked with numpy float16. Rule of 2,048: Model Training, \S8.6.3; Micikevicius et al. 2018.}
  \note{[3 min] [OPTIONAL] ENGAGE: students write both answers. Most write 0.5001 and 1.5. Step 2: FP16 numbers near 0.5 sit $2^{-11}$, about $4.9\times10^{-4}$, apart. The true sum, 0.5001 (open circle), is closer to 0.5 than to the next value, 0.50049, so it rounds back to 0.5. Step 3: repeat 10,000 times. FP32 climbs to 1.5; FP16 stays at 0.5. This is swamping. An update more than about 2,048 times smaller than the weight rounds away. Learning rate times gradient is usually that small, which is why mixed precision keeps a 32-bit master copy of the weights for the update. WARN: BF16 does not fix this; its spacing near 0.5 is $2^{-8}$, about $3.9\times10^{-3}$, coarser still. BF16 fixes range and leaves precision coarse. WARN: this is different from underflow; the update itself is representable, and the rounding of the sum loses it. FLEX: if cut, the loss-scaling slide states the underflow.}
\end{frame}

\begin{frame}{Loss scaling moves small gradients into FP16's range}
  \centering
  \mlsyssteps[width=0.94\textwidth,height=4.45cm,keepaspectratio]{loss_scale_step}{2}\\[0.0cm]
  {\normalsize\uncover<2->{Scaling the loss by $S$ adds $\log_2 S$ to every gradient's exponent.}}\\[0.05cm]
  {\footnotesize\uncover<3->{\textbf{Dynamic scaler:} overflow past 65,504 $\rightarrow$ skip the step, halve $S$.\enspace 2,000 clean steps $\rightarrow$ double $S$.}}
  \mlsyscredit{Data: Micikevicius et al., ICLR 2018, Figure 3 (SSD activation gradients, FP32 training; bins read from the figure, redrawn) and \S3.2 ($\times 8$). Rule: Model Training, \S8.6.3.}
  \note{[2.5 min] [CORE] Step 1: measured data. Micikevicius and colleagues logged the activation gradients of the SSD object detector during FP32 training; these are their bars, redrawn at slide size. 67 percent are exactly zero. Of the rest, about four in five (82 percent, read from the bars) lie below $2^{-24}$ and would become zero in FP16, while FP16's range above $2^{-10}$ goes unused. The paper reports that this network diverges in FP16 without scaling. Step 2: multiply the loss by $S$ before the backward pass. Differentiation is linear, so every gradient is multiplied by $S$, a shift of $\log_2 S$ on this axis; dividing by $S$ in 32 bits before the update loses nothing. The paper found $S = 8$, a shift of 3, enough to match FP32. The headroom is large, because the largest bin ends 26 doublings below 65,504. Step 3: a dynamic scaler searches for $S$. At $S = 1{,}024$ a gradient of $10^{-6}$ becomes $1.02\times10^{-3}$, and any gradient above 64 overflows; the scaler then skips the step and halves $S$. LINK: Lecture 7's GradScaler is this scaler. BF16 usually needs no scaler, because its floor is about $10^{-38}$.}
\end{frame}

\begin{frame}{Mixed precision keeps 32 bits where values accumulate}
  \centering
  \mlsyssteps[width=0.95\textwidth,height=4.75cm,keepaspectratio]{mixed_precision_step}{3}\\[0.1cm]
  {\normalsize\uncover<3->{Master copy and moments stay 32-bit, so a parameter still costs \textbf{16 B}.}}
  \mlsyscredit{Model Training, \S8.6.3 (interactive Explore figure: FP16 Mixed-Precision Training).}
  \note{[2 min] [CORE] The conventional FP16 recipe in six stages. Step 1: stages 1 and 2. The 32-bit master weights are cast to FP16, and the forward pass runs in FP16 on the Tensor Cores. Activations are stored in 16 bits, which is the memory saving. Step 2: stages 3 to 5. The loss is scaled, the backward pass produces scaled FP16 gradients, and they are copied to 32 bits and unscaled. Step 3: stage 6. The update is applied to the FP32 master weights, where the 0.0001 survives. The static state is unchanged from Part 1: 16 bytes per parameter, because the master copy and the moments stay in 32 bits. Mixed precision buys arithmetic rate and activation bytes and leaves the static bytes unchanged. ENGAGE: which stages does the BF16 recipe drop, and why? (Expected: stages 3 and 5, the loss scaling and unscaling, because BF16 has FP32's range; the master copy stays, because BF16 has even fewer mantissa bits than FP16.) LINK: Lecture 7's autocast keeps softmax and reductions in FP32.}
\end{frame}

\begin{frame}{Diagnose: what went wrong in each FP16 run?}
  \mlsysactivity{Diagnose}{2.5 min}
  \begin{columns}[T]
    \begin{column}{0.43\textwidth}
      \centering
      {\footnotesize\textbf{\textcolor{GreenLine}{Healthy: a real FP16 run}}}\\[0.05cm]
      \safeimg[width=\textwidth,height=4.5cm,keepaspectratio]{scaler_trace.pdf}
    \end{column}
    \begin{column}{0.55\textwidth}
      \small
      \textbf{\textcolor{BlueLine}{Run A.}} The scale halves every few dozen steps; many steps are skipped.\\[0.05cm]
      \uncover<2->{\textbf{Overflow:} scaled gradients pass 65,504. Look for a gradient spike or a bad batch.}\\[0.4cm]
      \textbf{\textcolor{OrangeLine}{Run B.}} The scale stays at 1. Loss flattens early; late-layer gradients read exactly zero.\\[0.05cm]
      \uncover<2->{\textbf{Underflow:} small gradients fall below FP16's floor. Enable dynamic scaling, or use BF16.}
    \end{column}
  \end{columns}
  \mlsyscredit{Left: our run, torch.amp.GradScaler at its defaults, two-layer MLP in FP16 autocast, PyTorch 2.13, Apple M5 Max GPU (images/scaler\_run.py). Model Training, \S8.6.3.}
  \note{[2.5 min] [OPTIONAL] Left: a real trace, produced by running PyTorch's GradScaler at its default settings on a small FP16 model (images/scaler\_run.py; 28,000 steps). The scale starts at $2^{16}$ and doubles after every 2,000 clean steps, nine times, to $2^{25}$. There the first overflow appears at step 18,001: the scaler skips that step and halves the scale, and the cycle repeats every 2,001 steps, 5 skipped steps in the run. A healthy scaler sits just under the overflow ceiling and skips a step rarely. ENGAGE: pairs name the failure in each run and one fix, two minutes. Step 2: Run A overflows. Constant halving means some gradients are genuinely huge, so look for a loss spike, a bad batch, or a learning rate that is too high. Run B underflows. With the scale fixed at 1, small late-layer gradients fall below FP16's normal floor of $6.1\times10^{-5}$ and become zero below about $6\times10^{-8}$; those weights stop learning. Dynamic loss scaling or BF16 fixes it. WARN: students blame the learning rate for Run B; a zero gradient has nothing to do with step size.}
\end{frame}

\begin{frame}{16-bit arithmetic: 518 days to 35, and 598 GB to 311}
  \begin{columns}[c]
    \begin{column}{0.46\textwidth}
      \centering\safeimg[width=\textwidth,height=4.2cm,keepaspectratio]{precision_days_col.pdf}
    \end{column}
    \begin{column}{0.52\textwidth}
      \centering\uncover<2->{\safeimg[width=\textwidth,height=4.2cm,keepaspectratio]{scoreboard_step2.pdf}}
    \end{column}
  \end{columns}
  \vspace{0.3cm}
  \centering{\normalsize 15$\times$ is a ceiling; elementwise, optimizer, and input work stay off the Tensor Cores.}
  \mlsyscredit{Model Training, Napkin Math 8.6 and 8.7 (50\% utilization held fixed, \$4 per hour).}
  \note{[1.5 min] [CORE] Left, run time for GPT-2 XL on one H100 at 50 percent utilization: FP32 about 518 days and \$49,751; TF32 about 70 days; BF16 about 35 days and \$3,370. Step 2: memory. Mixed precision halves the activations, 574 GB to 287 GB; the 24 GB of static state does not move. The step is 311 GB, still 3.6 H100s. WARN: holding utilization fixed overstates the gain. By Amdahl's law the work that does not run on Tensor Cores (normalization, elementwise operations, the optimizer update, input waits) bounds the real speedup below 15 times. Land the segment: 16-bit formats raise the arithmetic ceiling and halve the activations, but small updates and small gradients need 32-bit master weights and loss scaling, and the 16 bytes per parameter stay.}
\end{frame}

% =============================================================================
% SECTION: TILING
% =============================================================================
\mlsysnote{[0.5 min] [CORE] Most of the 311 GB is one tensor. Find it, then stop storing it.}
\mlsysdivider{Why does attention move more bytes than its arithmetic needs?}{Model Training, \S8.3.3, \S8.3.4, and \S8.6.4}
\section{Tiling}

\begin{frame}{In one GPT-2 XL layer, attention scores take 83\% of the bytes}
  \centering
  {\small\textcolor{midgray}{\textbf{From Part 1:} one layer's saved activations, batch 32 of 1,024 tokens.}}\\[0.1cm]
  \mlsyssteps[width=0.95\textwidth,height=5.3cm,keepaspectratio]{layer_bytes_step}{2}
  \mlsyscredit{Model Training, Napkin Math 8.3 and \S8.6.4. The split of the score bar follows its byte counts.}
  \note{[1.5 min] [CORE] Step 1, a 30-second recall from Part 1: one layer's saved activations at batch 32 in 16 bits. The attention scores are 4,194 MB of the layer's 5,043 MB, 83 percent, and across 48 layers about 201 of the 287 GB. Step 2: the score bar is three saved copies of every entry, the softmax probabilities P (2 bytes), a one-byte dropout mask (dropout randomly zeroes values during training), and P after dropout (2 bytes), 5 bytes per entry. FlashAttention stores none of the three. It keeps, per row, the running maximum and sum that the next slides build, and it recomputes the rest in the backward pass. That is the 201 GB the tiling segment removes. WARN: the per-layer line items sum to 242 GB; the library's 287 GB also counts saved inputs this list omits.}
\end{frame}

\begin{frame}{Standard attention crosses HBM four times; tiling stays on chip}
  \centering
  \mlsyssteps[width=0.92\textwidth,height=5.45cm,keepaspectratio]{attention_io_step}{4}
  \mlsyscredit{Model Training, \S8.3, \S8.3.4, \S8.6.4 (interactive Explore figure: Attention Across the Memory Boundary). 228 KB/SM: Hardware Acceleration.}
  \note{[2.5 min] [CORE] Two memory tiers on the accelerator. Step 1: HBM holds 80 GB at 3.35 TB/s. Step 2: on-chip SRAM holds 30 to 50 MB across the chip at 15 to 20 TB/s (228 KB per SM on H100). Step 3: standard attention runs separate kernels and bypasses SRAM: (1) write scores $S$ to HBM, (2) read $S$, (3) write probabilities $P$, (4) read $P$. Counting only scores, intensity is $d_{\text{head}}/2 = 32$ FLOP/byte against the H100's ridge of 295, so arithmetic stalls on HBM. Step 4: FlashAttention tiles attention inside SRAM. Query, key, and value blocks stream into SRAM; $S$ and $P$ never leave on-chip memory, online softmax maintains running $(m, \ell)$ on chip, and only the final output $O$ and $(m, \ell)$ cross back to HBM. LINK: Lecture 2 introduced the roofline; Lecture 11 develops the full hierarchy.}
\end{frame}

\begin{frame}{On GPT-2, mask, softmax, and dropout outweigh the matmuls}
  \centering
  \begin{columns}[c]
    \begin{column}{0.50\textwidth}
      \centering\safeimg[width=\textwidth,height=4.4cm,keepaspectratio]{fa_fig1_hierarchy.png}
    \end{column}
    \begin{column}{0.42\textwidth}
      \centering\uncover<2->{\safeimg[width=\textwidth,height=4.6cm,keepaspectratio]{fa_fig1_gpt2.png}}
    \end{column}
  \end{columns}
  \vspace{0.15cm}
  \uncover<2->{{\normalsize One fused kernel instead of five made attention 7.6$\times$ faster on an A100.}}
  \mlsyscredit{Dao et al., FlashAttention, NeurIPS 2022, Figure 1, left and right panels (NVIDIA A100), arxiv.org/abs/2205.14135.}
  \note{[1.5 min] [OPTIONAL] Two panels from Dao et al.'s Figure 1 on an NVIDIA A100. Step 1, left: the A100's two tiers, 20 MB of SRAM at 19 TB/s against 40 GB of HBM at 1.5 TB/s (12.7 times slower). Step 2, right: PyTorch GPT-2 attention runtime split by kernel. The two matmuls are the thin slices at the bottom and top; mask, softmax, and dropout do almost no FLOPs yet dominate the bar because each reads and writes the full score matrix in HBM. Fusing all five into one tiled SRAM kernel ran 7.6 times faster. ENGAGE: which slices would a faster matrix unit shrink? (Expected: only the two matmuls.) WARN: these are A100 numbers; the SRAM-to-HBM bandwidth gap carries over to H100. FLEX: cut freely.}
\end{frame}

\begin{frame}{A 128 $\times$ 128 tile is 65.5 KB; the full matrix is 67.1 MB}
  \centering
  \mlsyssteps[width=0.95\textwidth,height=4.7cm,keepaspectratio]{tile_grid_step}{2}\\[0.05cm]
  {\small\textcolor{midgray}{Chapter example: 4,096 tokens, one head, FP32 scores.}}
  \mlsyscredit{Model Training, \S8.6.4 (4,096 tokens, one head, FP32 scores); 228 KB per SM: Hardware Acceleration. Tile count derived.}
  \note{[2 min] [CORE] Step 1: one head's score matrix at 4,096 tokens, cut into tiles of 128 by 128, 32 by 32 tiles, 1,024 in all. FlashAttention takes one block of 128 queries (the blue row) and streams the key and value blocks past it, one tile at a time, entirely in on-chip memory. Step 2: the sizes. The full FP32 matrix is 67.1 MB per head; 16 heads would be 1.1 GB. One tile is 65.5 KB, which fits under the 228 KB of shared memory one SM can hold. ENGAGE: what is the obstacle? (Expected: softmax divides each score by a sum over the whole row, and one tile sees only part of the row.) The next slide removes that obstacle. WARN: the chapter counts FP32 scores here because softmax often runs in FP32; in FP16 every number halves and the quadratic growth stays.}
\end{frame}

\begin{frame}{Softmax can be computed one block at a time}
  \only<3|handout:0>{\mlsysactivity{Your turn}{2 min}}
  \centering
  \begin{tikzpicture}
    \node[anchor=east, font=\small\bfseries, text=darktext] at (-0.8,0) {scores $x$};
    \foreach \x/\v in {0/1, 1.2/3, 2.6/2, 3.8/5}{\scorecell{\x}{\v}{white}}
    \only<2->{\scorecell{0}{1}{BlueL}\scorecell{1.2}{3}{BlueL}}
    \only<4->{\scorecell{2.6}{2}{OrangeL}\scorecell{3.8}{5}{OrangeL}}
    \node[font=\footnotesize\bfseries, text=BlueLine] at (0.6,0.65) {block 1};
    \node[font=\footnotesize\bfseries, text=OrangeLine] at (3.2,0.65) {block 2};
    \foreach \x/\v in {0/10, 1.2/20, 2.6/30, 3.8/40}{%
      \node[font=\small\bfseries, text=GreenLine] at (\x,-0.62) {\v};}
    \node[anchor=east, font=\small\bfseries, text=GreenLine] at (-0.8,-0.62) {values $v$};
  \end{tikzpicture}\\[0.2cm]
  \small
  \begin{tabular}{@{}llll@{}}
    \toprule
    & max $m$ & sum $\ell = \sum e^{x - m}$ & output $O = \sum e^{x - m}\, v$ \\
    \midrule
    Whole row & 5 & $e^{-4} + e^{-2} + e^{-3} + 1 = \mathbf{1.203}$ & $\mathbf{44.38}$ \\
    \uncover<2->{Block 1} & \uncover<2->{3} & \uncover<2->{$e^{-2} + 1 = 1.135$} & \uncover<2->{$10e^{-2} + 20 = 21.35$} \\
    \uncover<4->{Block 2} & \uncover<4->{5} & \uncover<4->{$1.135\, e^{-2} + e^{-3} + 1 = \mathbf{1.203}$} & \uncover<5->{$21.35\, e^{-2} + 30e^{-3} + 40 = \mathbf{44.38}$} \\
    \bottomrule
  \end{tabular}\\[0.2cm]
  {\normalsize\only<3|handout:0>{\textcolor{mlsysaccent}{\textbf{Your turn:} block 2 holds 2 and 5. Finish $m$ and $\ell$.}}%
  \uncover<5->{Result: $O/\ell = 44.38 / 1.203 = \mathbf{36.9}$, the same as the whole row.}}
  \mlsyscredit{Toy example, derived. Algorithm: Model Training, \S8.6.4; Milakov and Gimelshein 2018; Dao et al. 2022.}
  \note{[4.5 min] [CORE] Four scores in two blocks, each with a value $v$; the row's output is the softmax-weighted average of the values. Step 1: direct whole-row result: $m = 5$, $\ell = 1.203$, $O = 44.38$. Step 2: block 1 alone gives running $m = 3$, $\ell = 1.135$, and unnormalized output $O = 10e^{-2} + 20 = 21.35$. Step 3: ENGAGE, your turn in pairs, two minutes. Hint: the old sum was scaled relative to 3; what happens when the maximum rises to 5? Step 4: rescale old $\ell$ by $e^{3-5} = 0.135$, giving 0.154, then add $e^{-3} = 0.050$ and 1 to recover 1.203 exactly. Step 5: the same factor $e^{3-5}$ rescales old $O$ ($21.35 \times 0.135 = 2.89$), plus $30e^{-3} = 1.49$ plus 40 = 44.38. Divide once at the end: $O/\ell = 36.9$, identical to the whole row. Per row, SRAM keeps only $m$, $\ell$, and running $O$. WARN: exact, not approximate.}
\end{frame}

\begin{frame}{Backward keeps two numbers per row in place of the matrix}
  \begin{columns}[c]
    \begin{column}{0.66\textwidth}
      \centering\safeimg[width=\textwidth,height=3.6cm,keepaspectratio]{flash_store.pdf}\\[0.2cm]
      {\normalsize 32.7$\times$ less stored per head; backward recomputes the tiles.}
    \end{column}
    \begin{column}{0.31\textwidth}
      \centering\safeimg[width=\textwidth,height=2.8cm,keepaspectratio]{amd-fiji-hbm.jpg}\\[0.08cm]
      {\footnotesize\textcolor{midgray}{Four HBM stacks flanking a GPU die: avoiding off-die writes saves both space and bandwidth.}}
    \end{column}
  \end{columns}
  \mlsyscredit{Model Training, \S8.6.4 (4,096 tokens, FP32). Photo: C.\ Spille/pcgameshardware.de (CC BY-SA 4.0), Wikimedia Commons.}
  \note{[1.5 min] [OPTIONAL] What each kernel keeps for the backward pass, per head at 4,096 tokens. Standard attention keeps its inputs, its output, and both the score and probability matrices, about 138.4 MB. FlashAttention keeps the queries, keys, values, the output, and the running maximum and sum for each row, about 4.2 MB, 32.7 times less. The gap widens with sequence length, because only the standard footprint is quadratic. During backward, FlashAttention recomputes each score tile from Q and K in on-chip memory instead of reading it from HBM. ENGAGE: recomputing costs arithmetic. Is that a good trade on this machine? Hold the answer for the next slide. FLEX: if short on time, state the 32.7 times and move on.}
\end{frame}

\begin{frame}{More FLOPs and fewer bytes run 5.7$\times$ faster}
  \mlsysactivity{Predict}{2 min}
  {\normalsize FlashAttention recomputes tiles, so it does more arithmetic. Faster or slower?}\\[0.25cm]
  \centering
  \uncover<2->{%
    \begin{tabular}{@{}lrr@{}}
      \toprule
      & Standard & FlashAttention \\
      \midrule
      GFLOPs & 66.6 & 75.2 \\
      HBM reads and writes & 40.3 GB & \textbf{4.4 GB} \\
      Runtime & 41.7 ms & \textbf{7.3 ms} \\
      \bottomrule
    \end{tabular}\\[0.1cm]
    {\footnotesize\textcolor{midgray}{GPT-2 medium attention, forward and backward, 1,024 tokens, batch 64, A100.}}}\\[0.35cm]
  {\Large\uncover<3->{GPT-2 XL step: 311 GB $\rightarrow$ \textbf{109.6 GB}}}
  \mlsyscredit{Table: Dao et al. 2022, Figure 2 (left). GPT-2 XL step: Model Training, \S8.6.4 and Napkin Math 8.10.}
  \note{[2.5 min] [CORE] ENGAGE: vote faster or slower, with one reason. Step 2: measured data from the FlashAttention paper. It does about 13 percent more arithmetic, 75.2 against 66.6 GFLOPs, because backward recomputes the tiles. It moves about 9 times fewer bytes through HBM, 4.4 against 40.3 GB, and runs about 5.7 times faster, 7.3 against 41.7 ms. Bytes limited the kernel, so trading FLOPs for bytes wins. Step 3: back to the running case. Without the stored scores, GPT-2 XL's step drops from 311 GB to 109.6 GB, still about 24 GB over the H100's 85.9 GB; the scoreboard returns at the end of the next segment.}
\end{frame}

\begin{frame}{You already call it, because one PyTorch function picks the kernel}
  \centering\vspace{-0.18cm}
  {\small\ttfamily out = F.scaled\_dot\_product\_attention(q, k, v, is\_causal=True)}\\[0.04cm]
  {\footnotesize\textcolor{midgray}{GPT-2 XL attention, one sequence: 25 heads, 1,024 tokens, head width 64, FP32, CPU.}}\\[0.12cm]
  {\footnotesize\ttfamily
    \begin{tabular}{@{}lrr@{}}
      \multicolumn{3}{@{}l}{\normalfont\bfseries\color{GreenLine}== default dispatch} \\
      operator & self MB & total ms \\
        aten::\_scaled\_dot\_product\_flash\_attention\_for\_cpu & -3.6 & \textcolor{GreenLine}{\textbf{13.16}} \\
        aten::empty\_strided & \textcolor{GreenLine}{\textbf{6.6}} & 0.17 \\
      \multicolumn{3}{@{}l}{} \\[-0.25cm]
      \multicolumn{3}{@{}l}{\normalfont\bfseries\color{RedLine}== forced math backend} \\
        aten::\_scaled\_dot\_product\_attention\_math & -124.3 & \textcolor{RedLine}{\textbf{45.23}} \\
        aten::bmm & \textcolor{RedLine}{\textbf{111.4}} & 21.13 \\
        aten::\_safe\_softmax & -26.2 & 19.02 \\
    \end{tabular}}\\[0.12cm]
  {\normalsize\uncover<2->{Math path: a 104.9 MB score matrix and 3.4$\times$ the time; same output to $10^{-6}$.}}
  \mlsyscredit{Our torch.profiler run, PyTorch 2.13, Apple M5 Max CPU (images/sdpa\_profile.py; rows from images/data/sdpa\_profile.txt). API: PyTorch docs, scaled\_dot\_product\_attention.}
  \note{[1 min] [OPTIONAL] Students do not write FlashAttention; they call it. This line runs GPT-2 XL attention for one sequence under torch.profiler on a laptop CPU (images/sdpa\_profile.py). By default it dispatches to the fused flash kernel, allocating only the 6.6 MB output in 13.16 ms. Step 2: forced to the reference math backend, batched matmul allocates 111.4 MB (the 104.9 MB score matrix plus output) and runs in 45.23 ms, 3.4 times slower, with outputs agreeing to $10^{-6}$. On a GPU in FP16 or BF16, the same call dispatches to FlashAttention CUDA kernels. WARN: negative self MB means an op freed an intermediate tensor; CPU times show dispatch, not GPU speed. Land the segment: FlashAttention tiles exact attention on chip, keeping two numbers per row instead of the score matrix. FLEX: cut freely.}
\end{frame}

% =============================================================================
% SECTION: TRADE-OFFS
% =============================================================================
\mlsysnote{[0.5 min] [CORE] The rest is ordinary activations. Now something must be given up.}
\mlsysdivider{When the activations still do not fit, what do we give up?}{Model Training, \S8.6.5 and \S8.6.6}
\section{Trade-offs}

\begin{frame}{After FlashAttention, 23 of the 32 sequences fit}
  \centering
  \mlsyssteps[width=0.92\textwidth,height=5.0cm,keepaspectratio]{batch_fit_step}{2}
  \mlsyscredit{Derived from Model Training, Napkin Math 8.10 (85.6 GB of activations at batch 32; 24 GB static).}
  \note{[2 min] [CORE] Step memory against the number of sequences resident at once. The line starts at 24 GB of static state and rises by about 2.67 GB per sequence, 85.6 GB of activations divided by 32. It crosses the H100's 85.9 GB near 23 sequences. Step 2: batch 16 needs 66.8 GB and fits; batch 32 needs 109.6 GB and does not. ENGAGE: why not train at batch 23? (Expected: possible, but the recipe and its learning rate were tuned for 32, and changing the batch changes the optimization.) Two ways to keep the batch of 32 follow.}
\end{frame}

\begin{frame}{Accumulation adds micro-batch gradients, then updates once}
  \centering
  \mlsyssteps[width=0.95\textwidth,height=3.65cm,keepaspectratio]{grad_accum_step}{3}\\[0.15cm]
  {\normalsize $\displaystyle \frac{1}{B}\sum_{i=1}^{B}\nabla \ell_i \;=\; \frac{1}{K}\sum_{j=1}^{K}\Bigl(\frac{1}{b}\sum_{i\in\text{micro-batch } j}\nabla \ell_i\Bigr), \qquad B = K\,b$}\\[0.12cm]
  {\normalsize\uncover<3->{$2 \times 16$: \textbf{66.8 GB}, same FLOPs, one update per two passes.}}
  \mlsyscredit{Model Training, \S8.6.5, Eq.\ 10, Napkin Math 8.10 (interactive Explore figure: Gradient Accumulation).}
  \note{[2.5 min] [CORE] Step 1: run a micro-batch forward and compute its loss. Step 2: run backward; its gradient is added into the gradient buffer, and the weights do not change. Step 3: after $K$ micro-batches, the summed gradient, divided by $K$ (the orange tag), drives one optimizer step. The equation says why this is the same gradient. Averaging $K$ micro-batch means of $b$ examples equals the mean over all $B = Kb$ examples. For GPT-2 XL, two micro-batches of 16: only 16 sequences' activations are alive at a time, so the step needs 66.8 GB, and the arithmetic per update is unchanged. The cost: one update per two passes, and smaller micro-batches can lower utilization. WARN: equivalence needs the same examples and consistent averaging; BatchNorm, which computes statistics per micro-batch, breaks it. GPT-2 uses LayerNorm, so it holds up to dropout and rounding order.}
\end{frame}

\begin{frame}{Diagnose: why are these updates 4$\times$ too large?}
  \mlsysactivity{Diagnose}{2.5 min}
  \begin{columns}[c]
    \begin{column}{0.62\textwidth}
      \footnotesize\ttfamily
      \begin{tabular}{@{}l@{}}
        opt = SGD(params, lr=0.1, momentum=0.9)\\
        for step, (x, y) in enumerate(loader):\\
        \hspace*{1.2em}loss = loss\_fn(model(x), y)\\
        \hspace*{1.2em}\alt<2->{\textcolor{GreenLine}{\textbf{(loss / 4).backward()}}}{\textcolor{RedLine}{\textbf{loss.backward()}}}\\
        \hspace*{1.2em}if (step + 1) \% 4 == 0:\\
        \hspace*{2.4em}opt.step()\\
        \hspace*{2.4em}opt.zero\_grad()
      \end{tabular}
    \end{column}
    \begin{column}{0.36\textwidth}
      \small
      Four micro-batches per update. Training diverges as if the learning rate were too high.\\[0.3cm]
      \uncover<2->{\textbf{Sum of four means} $=$ 4$\times$ the batch mean. Divide each loss by 4.}
    \end{column}
  \end{columns}
  \mlsyscredit{Adapted from Model Training, listing 8 and \S8.6.5 (practical considerations).}
  \note{[2.5 min] [OPTIONAL] ENGAGE: pairs find the bug, two minutes. The loss function returns the mean over its micro-batch, and backward adds into the gradient buffer. Step 2: after four micro-batches the buffer holds the sum of four means, four times the mean over all their examples, so every update is four times too large. Dividing each loss by 4, or dividing the learning rate by 4, restores the batch-mean convention. WARN: with Adam the bug hides, because Adam divides by a running gradient magnitude and the factor mostly cancels; gradient clipping at a norm of 1.0 and Adam's epsilon still change, so the run differs in ways that are hard to trace. The chapter calls loss normalization a common source of subtle bugs.}
\end{frame}

\begin{frame}{Checkpointing stores some activations and recomputes the rest}
  \centering
  \mlsyssteps[width=0.95\textwidth,height=4.85cm,keepaspectratio]{act_ckpt_step}{3}
  \mlsyscredit{Model Training, \S8.6.5 (interactive Explore figure: Activation Checkpointing).}
  \note{[1.5 min] [CORE] Activation checkpointing (distinct from disk checkpoints). Step 1: in the forward pass, only checkpoint boundaries ($a_0$ and $a_3$, green) stay in memory; interior activations ($a_1$, $a_2$, $a_4$, dashed gray) are dropped as soon as the next layer consumes them. Step 2: when backward reaches $a_4$, it replays forward from $a_3$ to recompute $a_4$ (orange), computes its gradient, and frees $a_4$. Step 3: backward then replays forward from $a_0$ to recompute $a_1$ and $a_2$, freeing each after use. ENGAGE: what does backward pay for every orange box? (Expected: one extra forward computation of that layer.) LINK: Lecture 7 showed torch.utils.checkpoint; the next slide sizes the segments for GPT-2 XL.}
\end{frame}

\begin{frame}{For GPT-2 XL, eight segment inputs replace 85.6 GB}
  \centering
  \vspace*{0.2cm}
  \mlsyssteps[width=0.92\textwidth,height=4.3cm,keepaspectratio]{ckpt_segments_step}{2}\\[0.2cm]
  {\normalsize\uncover<3->{Activations 85.6 $\rightarrow$ \textbf{11.5 GB}; step \textbf{35.5 GB}; one extra forward per layer.}}
  \mlsyscredit{Model Training, Napkin Math 8.8 (48 layers in 8 segments of 6) and \S8.6.5.}
  \note{[2 min] [CORE] Step 1: GPT-2 XL's 48 layers as 8 segments of 6. The forward pass keeps each segment's input (green bars), 32 by 1,024 by 1,600 values in 16 bits, 0.84 GB for all eight, and drops everything inside the segments. Step 2: the backward pass reruns one segment forward from its input (orange), uses those activations, frees them, and moves to the previous segment, with 10.7 GB live at a time. Step 3: activations fall from 85.6 GB to 11.5 GB, 7.4 times less, and the whole step to 35.5 GB. The cost: every layer runs forward twice, so a step that cost about three forward passes now costs four, about one-third more arithmetic.}
\end{frame}

\begin{frame}{Checkpoints every $\sqrt{L}$ layers minimize memory}
  \centering
  \mlsyssteps[width=0.92\textwidth,height=4.6cm,keepaspectratio]{sqrt_ckpt_step}{2}\\[0.1cm]
  {\small\uncover<2->{Measured: a 1,000-layer ResNet, 48 GB $\rightarrow$ 7 GB for 30\% more time (Chen et al. 2016).}}
  \mlsyscredit{Model Training, \S8.6.5.2 and table 13 (curve derived, equal layers); Chen et al. 2016, abstract.}
  \note{[2 min] [OPTIONAL] Activation memory in units of one layer, for $L = 48$ layers and $k$ evenly spaced checkpoints: $k$ stored checkpoints plus one segment of $48/k$ layers rebuilt at a time, the $k + 48/k$ printed on the curve. With one checkpoint the rebuilt segment is the whole network, 49 units, no saving. With 48 checkpoints every layer is stored, 49 units again. The minimum is at $k = \sqrt{48} \approx 7$: about 14 units. Step 2: 71 percent less activation memory for about one extra forward pass, 33 percent more step arithmetic. GPT-2 XL's 8 segments sit at $k = 8$, also 14 units. Chen and colleagues measured the trade on a 1,000-layer residual network, 48 GB to 7 GB for 30 percent more running time. ENGAGE: why is the extra cost about one forward pass whatever $k$ is? (Expected: each non-checkpoint layer is recomputed once.)}
\end{frame}

\begin{frame}{Recomputing only the attention core costs 7\% against 39\%}
  \centering
  \safeimg[width=0.95\textwidth,height=3.9cm,keepaspectratio]{selective_recompute.pdf}\\[0.2cm]
  {\normalsize Bytes saved per layer: $sbh\,(34 + 5as/h)$.\enspace GPT-2 XL: $34 + 80$.}\\[0.1cm]
  {\normalsize\uncover<2->{Dropping and recomputing the 80 (the scores) takes 287 $\rightarrow$ \textbf{85.6 GB}, 70.2\% less.}}
  \mlsyscredit{Korthikanti et al., Reducing Activation Recomputation in Large Transformer Models, MLSys 2023, arxiv.org/abs/2205.05198: Table 4 (one layer of a 22B GPT, A100) and Table 2. GPT-2 XL values derived.}
  \note{[2 min] [OPTIONAL] Checkpointing whole layers recomputes everything. Korthikanti and colleagues at NVIDIA measured a finer choice. Their Table 2 counts the bytes one transformer layer saves for backward: $sbh(34 + 5as/h)$, for sequence length $s$, batch $b$, hidden width $h$, and $a$ heads. The $5as/h$ term is the attention core, the score, softmax, and dropout tensors at 5 bytes per score entry, the same count as the layer-bytes slide. The bars are their measurement for one layer of a 22-billion-parameter GPT, with no recomputation at 19.6 ms, full recomputation 27.2 ms (39 percent more), recomputing only the attention core 20.9 ms (7 percent more). The core holds many bytes and few FLOPs. Step 2: for GPT-2 XL the core is 80 against 34, so the formula over 48 layers gives the chapter's 287 GB, and dropping the core leaves 85.6 GB, 70.2 percent less, the same bytes FlashAttention never stores. At scale, adding sequence parallelism, they trained a 530-billion-parameter model on 2,240 A100s at 54.2 percent model FLOPs utilization against 42.1 with full recomputation, 29 percent faster. ENGAGE: why recompute the attention core and not the feed-forward layers? (Expected: most bytes for the fewest FLOPs.)}
\end{frame}

\begin{frame}{Vote: accumulate, checkpoint, or both?}
  \mlsysactivity{Vote}{2.5 min}
  \begin{columns}[T]
    \begin{column}{0.40\textwidth}
      \small
      \choice{A}{Accumulate two micro-batches of 16}\\[0.1cm]
      {\small\uncover<2->{\textcolor{GreenLine}{66.8 GB; same FLOPs; half as many updates.}}}\\[0.3cm]
      \choice{B}{Checkpoint eight segments}\\[0.1cm]
      {\small\uncover<2->{\textcolor{OrangeLine}{35.5 GB; one-third more FLOPs; 47 days.}}}\\[0.3cm]
      \choice{C}{Both}
    \end{column}
    \begin{column}{0.58\textwidth}
      \centering
      \uncover<3->{\safeimg[width=\textwidth,height=5.4cm,keepaspectratio]{scoreboard_step4.pdf}}
    \end{column}
  \end{columns}
  \mlsyscredit{Model Training, Napkin Math 8.8 and 8.10.}
  \note{[2.5 min] [CORE] ENGAGE: vote A, B, or C, with one cost for your choice. Step 2: both fit (the numbers appear on the cards). Accumulation needs 66.8 GB and keeps the arithmetic, so the run stays about 35 days; it pays in update frequency, and 16-sequence micro-batches may use the chip less well. Checkpointing needs only 35.5 GB but adds one-third more arithmetic, about 47 days at the same rate. The chapter takes A, because the batch it needs fits without recomputation. B is the choice when a larger resident micro-batch or a deeper model matters more than arithmetic; C when neither alone fits. Step 3: the scoreboard. 598 to 66.8 GB, nine times smaller, on one H100; the dashed orange bar beside it is option B, checkpointing, at 35.5 GB. The opening vote's answer was D. Land the segment: accumulation gives up update frequency, checkpointing gives up arithmetic, and either fits the step.}
\end{frame}

\begin{frame}{In Hugging Face, today's techniques are four arguments}
  \begin{columns}[T]
    \begin{column}{0.48\textwidth}
      {\footnotesize\ttfamily
      args = TrainingArguments(\\
      \hspace*{1.2em}per\_device\_train\_batch\_size=16,\\
      \hspace*{1.2em}gradient\_accumulation\_steps=2,\\
      \hspace*{1.2em}gradient\_checkpointing=True,\\
      \hspace*{1.2em}bf16=True,\\
      )\\[0.25cm]
      from\_pretrained(...,\\
      \hspace*{1.2em}attn\_implementation="sdpa")}
    \end{column}
    \begin{column}{0.50\textwidth}
      \footnotesize
      \renewcommand{\arraystretch}{1.3}
      \begin{tabular}{@{}lll@{}}
        \toprule
        \textbf{Argument} & \textbf{GPT-2 XL} & \textbf{Cost} \\
        \midrule
        bf16 & 598 $\rightarrow$ 311 GB & mantissa bits \\
        sdpa & $\rightarrow$ 109.6 GB & recomputed tiles \\
        accumulation & $\rightarrow$ 66.8 GB & fewer updates \\
        checkpointing & or 35.5 GB & $+\frac{1}{3}$ arithmetic \\
        \bottomrule
      \end{tabular}\\[0.25cm]
      {\textcolor{midgray}{Docs, on checkpointing: ``slower training speed (\textasciitilde 20\%)''.}}
    \end{column}
  \end{columns}
  \mlsyscredit{Code, quote: Hugging Face Transformers v5.0.0, huggingface.co/docs/transformers/perf\_train\_gpu\_one (Apache-2.0); batch values set to the running case (guide: 4 and 16). GPT-2 XL: Napkin Math 8.8, 8.10.}
  \note{[2 min] [CORE] This is the code students will write. The TrainingArguments block is the Hugging Face guide's for efficient training on one GPU, with its two batch values set to today's running case; the same page loads a model with attn\_implementation="sdpa", which selects the fused attention of the last segment (the model name and device argument are omitted here). Walk the table. Each argument is one segment of today, with its GPT-2 XL effect and its bill. bf16 halves the activations and keeps FP32 master weights in the optimizer. sdpa stops storing the scores. gradient\_accumulation\_steps keeps the effective batch with one micro-batch resident: on the slide, 16 sequences twice for GPT-2 XL's batch of 32 (the docs' example runs 4 sequences 16 times for a batch of 64). gradient\_checkpointing reaches 35.5 GB instead; the docs measure about 20 percent slower training, while one-third more is our count of arithmetic and was never timed. WARN: the docs advise against many accumulation steps on tiny micro-batches (they mark batch 1 times 64 down and 4 times 16 up), because small micro-batches underuse the accelerator. ENGAGE: if the step already fits, which argument would you remove first? (Expected: gradient checkpointing, the only one that adds arithmetic.)}
\end{frame}

\begin{frame}{Beyond one GPU: DDP replicates state; ZeRO/FSDP shards it}
  \begin{columns}[c]
    \begin{column}{0.67\textwidth}
      \centering\mlsyssteps[width=\textwidth,height=4.45cm,keepaspectratio]{zero_shard_step}{3}
    \end{column}
    \begin{column}{0.30\textwidth}
      {\centering\safeimg[width=\textwidth,height=2.2cm,keepaspectratio]{nvidia-dgx-b200-hgx.jpg}\par}\vspace{0.12cm}
      \raggedright\footnotesize
      \textbf{8$\times$ H100 NVLink} (450 GB/s/dir):\\[0.08cm]
      GPT-2 XL Ring AllReduce: \textbf{12 ms} (2\% of 0.60 s step).\\[0.08cm]
      \uncover<3->{\textcolor{mlsysaccent}{\textbf{ZeRO-3 / FSDP:}} 8$\times$ less state per GPU at \textbf{1.5$\times$} comm. volume.}
    \end{column}
  \end{columns}
  \mlsyscredit{Model Training, \S8.7 (8$\times$ H100 NVLink; Rajbhandari et al.\ SC 2020). Photo: Pokiiri (CC BY-SA 4.0), Wikimedia Commons.}
  \note{[2.5 min] [CORE] When static state alone outgrows one GPU (like Part 1's 128 GB Llama 3 8B full fine-tune), single-node scaling splits work across an 8-GPU NVLink baseboard (right). Step 1: Distributed Data Parallel (DDP, ZeRO-0) replicates all 16 B/param on every GPU and averages gradients with Ring AllReduce. For GPT-2 XL's 3 GB of gradients over 450 GB/s NVLink, Ring AllReduce moves 5.25 GB in about 12 ms (2 percent of a 0.60 s step), but Llama 3 8B still overflows at 128 GB per GPU. Step 2: ZeRO-1 shards the 12 B optimizer state (128 to 44 GB); ZeRO-2 also shards gradients (to 30 GB), with the same communication volume as DDP. Step 3: ZeRO-3 (PyTorch FSDP) shards weights too, cutting static state 8-fold to 16 GB per GPU (3 GB for GPT-2 XL) so full fine-tuning fits on one node, at 1.5 times the collective traffic. ENGAGE: which stage removes the biggest slice first? (Expected: ZeRO-1, optimizer state.)}
\end{frame}

% =============================================================================
% SECTION: WRAP-UP
% =============================================================================
\mlsysnote{[0.5 min] [CORE] Check the ideas, then return to AlexNet.}
\mlsysdivider{What did fitting the step cost?}{Model Training, \S8.6.6 and \S8.8}
\section{Wrap-Up}

\begin{frame}{Vote: true or false?}
  \mlsysactivity{Vote}{2.5 min}
  \small
  \renewcommand{\arraystretch}{1.35}
  \begin{tabular}{@{}>{\raggedright\arraybackslash}m{0.48\textwidth}>{\raggedright\arraybackslash}m{0.46\textwidth}@{}}
    \toprule
    \textbf{Claim} & \uncover<2->{\textbf{Verdict}} \\
    \midrule
    FlashAttention reduces attention's FLOPs. & \uncover<2->{\textbf{\textcolor{RedLine}{False.}} Slightly more FLOPs, far fewer bytes.} \\
    Accumulating $K$ micro-batches multiplies activation memory by $K$. & \uncover<2->{\textbf{\textcolor{RedLine}{False.}} One micro-batch is resident at a time.} \\
    BF16 needs no FP32 master weights. & \uncover<2->{\textbf{\textcolor{RedLine}{False.}} BF16 fixes range and leaves precision coarse.} \\
    Mixed precision makes training 15$\times$ faster on an H100. & \uncover<2->{\textbf{\textcolor{RedLine}{False.}} 15$\times$ is the peak ratio; Amdahl bounds it.} \\
    \bottomrule
  \end{tabular}
  \mlsyscredit{Model Training, \S8.6 self-check and \S8.8.}
  \note{[2.5 min] [CORE] ENGAGE: students write T or F for each, then compare with a neighbor. Step 2: all four are false. One: FlashAttention does about 13 percent more arithmetic and moves about 9 times fewer bytes. Two: only one micro-batch's activations are alive at a time; the chapter's self-check asks exactly this. Three: BF16's spacing near 0.5 is about $3.9\times10^{-3}$, so small updates vanish without a 32-bit master copy. Four: the 15 times is the ratio of peak rates; normalization, elementwise work, the optimizer, and input waits do not run on Tensor Cores. If many students marked the third true, return to the swamping slide.}
\end{frame}

\begin{frame}{Back to AlexNet, where one H100 now holds the step at a price}
  \begin{columns}[c]
    \begin{column}{0.64\textwidth}
      \centering\safeimg[width=\textwidth,height=4.7cm,keepaspectratio]{scoreboard_costs.pdf}
    \end{column}
    \begin{column}{0.33\textwidth}
      \centering
      \uncover<2->{\mlsysphoto[width=0.75\textwidth,height=2.4cm,keepaspectratio]{nvidia-gtx580-die.jpg}\\[0.1cm]
      {\small 2012: AlexNet split across\\two 3 GB GTX 580s.\\[0.1cm]2022: one 80 GB H100,\\after three techniques.}}
    \end{column}
  \end{columns}
  \vspace{0.25cm}
  \centering{\small\uncover<3->{It fits; a profile now finds the binding term, a one-core tokenizer (48\% busy).}}
  \mlsyscredit{Model Training, Napkin Math 8.10 and \S8.6.6. Photo: Fritzchens Fritz (CC0), Wikimedia Commons.}
  \note{[3 min] [CORE] Step 1: the whole scoreboard with each bill. Mixed precision took 598 to 311 GB and cost a 32-bit master copy and loss scaling. FlashAttention took it to 109.6 GB and costs recomputed tiles in backward. Two micro-batches of 16 took it to 66.8 GB and cost one update per two passes; checkpointing instead would have cost 12 more days. The 16 bytes per parameter never moved. Step 2: AlexNet split its network across two 3 GB cards. GPT-2 XL's step is about 200 times larger, and three techniques made it fit without a second card. Step 3: fitting comes first because a step that does not fit gives no profile; then profile it to find the binding term. The first profile of the fitted run shows the accelerator busy about 48 percent of the time, waiting on single-core tokenization, the input problem of Lecture 4b and Part 1.}
\end{frame}

\begin{frame}{Takeaways}
  \normalsize
  \begin{itemize}\setlength{\itemsep}{0.25cm}
    \item 16-bit formats raise the arithmetic peak 15$\times$ and halve activations; 32-bit master weights and loss scaling keep the numerics.
    \item FlashAttention computes exact attention tile by tile, keeping two numbers per row instead of the score matrix.
    \item Accumulation trades update frequency for memory; checkpointing trades one forward pass.
    \item Fit the step first, then profile it; each technique costs numerics, arithmetic, or update frequency.
  \end{itemize}
  \note{[1 min] [CORE] One line per segment, plus the order of work. Ask students to star the one they would have gotten wrong before today.}
\end{frame}

\begin{frame}{Open questions}
  \normalsize
  \begin{itemize}\setlength{\itemsep}{0.35cm}
    \item FP8 halves the bytes again but needs per-tensor scaling. How low can training precision go before the scaling machinery costs more than it saves?
    \item Checkpointing costs GPT-2 XL 12 extra days; a second H100 costs a second rental. Which would you buy, and what numbers decide it?
    \item Which other kernels in a training step deserve FlashAttention's rewrite?
  \end{itemize}
  \note{[1 min] [OPTIONAL] No single answer. FLEX: use one as a two-minute discussion if time remains, or as a written assignment. For the second, push students to name the numbers, rental price per hour, the deadline, and whether two accelerators would scale well (\S8.7). LINK: Lecture 10b develops formats below 16 bits.}
\end{frame}

\begin{frame}{Read and explore}
  \footnotesize
  \begin{itemize}\setlength{\itemsep}{0.04cm}
    \item \textbf{Read:} Model Training, \S8.6.3 to \S8.7; redo Napkin Math 8.8 for 12 segments of 4.
    \item \textbf{Explore:} \textit{FP16 Mixed-Precision Training} (\S8.6.3), \textit{Attention Across the Memory Boundary} (\S8.6.4), \textit{Gradient Accumulation} and \textit{Activation Checkpointing} (\S8.6.5) on mlsysbook.ai/vol1/.
    \item Micikevicius et al., ``Mixed Precision Training,'' ICLR 2018; Kalamkar et al., ``A Study of BFLOAT16,'' 2019.
    \item Dao et al., ``FlashAttention: Fast and Memory-Efficient Exact Attention with IO-Awareness,'' NeurIPS 2022.
    \item Chen et al., ``Training Deep Nets with Sublinear Memory Cost,'' 2016; Korthikanti et al., MLSys 2023.
    \item Rajbhandari et al., ``ZeRO,'' SC 2020; H.\ He, ``Making Deep Learning Go Brrrr From First Principles,'' 2022.
  \end{itemize}
  \note{[0.5 min] [CORE] The chapter first; its walkthrough repeats today's scoreboard with every number. Each paper backs one idea: master weights and loss scaling (Micikevicius), why BF16 rarely needs scaling (Kalamkar), tiling and online softmax (Dao), the square-root schedule (Chen), and selective recomputation (Korthikanti). He's essay is a short outside explanation of compute, memory traffic, and overhead.}
\end{frame}

\begin{frame}[c]
  \centering
  {\small\textcolor{lightgray}{NEXT LECTURE \textbar{} DATA SELECTION}}\\[0.5cm]
  {\LARGE\bfseries Which data is worth training on?}\\[0.4cm]
  {\normalsize\textcolor{midgray}{Memory gated the training-time law; today opened the gate. Its numerator, $6PD$, never changed.}}
  \note{[0.5 min] [CORE] Close on the training-time law. Memory sits outside the law as a feasibility gate, because a step that does not fit has no run time at all. Over two lectures we made each FLOP cheaper and opened that gate; the work itself, $6PD$, never changed. Lecture 9 asks which of the $D$ tokens are worth a step.}
\end{frame}

\end{document}
